\section{Definitions}
\label{sec:definitions}

This section introduces the formal notions required to reason about
decidability of enterprise states under incomplete observability.
All definitions are expressed strictly in terms of the Ontology of
Continua (OC Core) and the executable specification
ETS.K\_EA.v1.1. No new primitives, operators, or interpretative layers
are introduced.

\subsection{Observable domain}
\label{subsec:observable-domain}

Let $K_{EA}$ denote an enterprise continuum with internal state space
$\Omega_{EA}$ and boundary $\partial\Omega_{EA}$, equipped with axes
$A_{EA}$, potentials $P_{EA}$, thresholds $\Theta_{EA}$, flows $J_{EA}$,
cycles $C_{EA}$, continuity measure $k_{EA}$, and internal time
$\tau_{EA}$, as defined in ETS.K\_EA.v1.1.

Observability is defined exclusively via the admissible observation
constructs already present in ETS. Let
\[
\Omega_{\mathrm{obs}} \subseteq \Omega_{EA}
\]
denote the observable projection of the enterprise continuum, and let
\[
\Sigma_{\mathrm{obs}}
\]
denote the space of admissible observable signatures. The observation
operator is defined as
\[
\Pi : \Omega_{EA} \to \Sigma_{\mathrm{obs}}.
\]

No assumptions are made about injectivity, surjectivity, continuity, or
invertibility of $\Pi$ beyond what is explicitly fixed in ETS.

\subsection{Observational equivalence}
\label{subsec:observational-equivalence}

\begin{definition}[Observational equivalence]
Two enterprise continua $K_{EA}^{(1)}$ and $K_{EA}^{(2)}$ are said to be
\emph{observationally equivalent} if their induced observable signature
sets coincide:
\[
\Pi\!\left(\Omega_{EA}^{(1)}\right)
=
\Pi\!\left(\Omega_{EA}^{(2)}\right)
\subseteq \Sigma_{\mathrm{obs}}.
\]
\end{definition}

Observational equivalence induces equivalence classes over the space of
enterprise continua. All members of such a class are indistinguishable
by any admissible observation, even if their internal structures differ
in axes, potentials, flows, cycles, or threshold configurations.

\subsection{Enterprise state predicates}
\label{subsec:state-predicates}

Enterprise states are not identified with full internal configurations
in $\Omega_{EA}$, but with \emph{state predicates} defined over
admissible OC/ETS elements.

Typical examples include predicates corresponding to admissible phase
outcomes (OK, INERTIA, COLLAPSE, STOP), satisfaction or violation of
threshold conditions, or membership in specific regions of
$\Omega_{EA}$ or $\partial\Omega_{EA}$.

\begin{definition}[State predicate]
A state predicate is a mapping
\[
\varphi : \Omega_{EA} \to \{\text{true},\text{false}\},
\]
where $\varphi$ is defined solely using OC/ETS-admissible constructs and
relations.
\end{definition}

No operational, temporal, inferential, or methodological interpretation
is attached to state predicates.

\subsection{Decidability under incomplete observability}
\label{subsec:decidability}

\begin{definition}[Decidability]
A state predicate $\varphi$ is said to be \emph{decidable under
incomplete observability} if, for all enterprise continua
$K_{EA}^{(1)}$ and $K_{EA}^{(2)}$ that are observationally equivalent,
the predicate agrees:
\[
\Pi\!\left(\Omega_{EA}^{(1)}\right)
=
\Pi\!\left(\Omega_{EA}^{(2)}\right)
\;\Rightarrow\;
\varphi\!\left(\Omega_{EA}^{(1)}\right)
=
\varphi\!\left(\Omega_{EA}^{(2)}\right).
\]
\end{definition}

Decidability is a purely logical property: it requires invariance of the
predicate across all members of an observational equivalence class. No
algorithmic, statistical, inferential, or operational interpretation is
implied.

\begin{definition}[Undecidability]
A state predicate $\varphi$ is \emph{undecidable under incomplete
observability} if there exist two observationally equivalent enterprise
continua $K_{EA}^{(1)}$ and $K_{EA}^{(2)}$ such that
\[
\varphi\!\left(\Omega_{EA}^{(1)}\right)
\neq
\varphi\!\left(\Omega_{EA}^{(2)}\right).
\]
\end{definition}

Undecidability expresses the logical impossibility of determining the
truth value of $\varphi$ from any admissible observation, regardless of
data volume, observational persistence, or analytical sophistication.

\subsection{STOP-only predicates}
\label{subsec:stop-only}

Certain predicates admit a one-sided form of certification tied to
boundary conditions.

\begin{definition}[STOP-only decidability]
A predicate $\varphi$ is \emph{STOP-only decidable} if:
\begin{itemize}
\item whenever $\varphi$ evaluates to STOP for an enterprise continuum,
      this fact is invariant across its entire observational
      equivalence class;
\item for all non-STOP outcomes, there exist observationally equivalent
      continua with differing predicate values.
\end{itemize}
\end{definition}

STOP-only decidability captures the diagnosability boundary emphasized
in ETS: certain terminal or boundary conditions can be certified, while
all non-terminal states remain logically undecidable under incomplete
observability.

\subsection{Separation of observability and decidability}
\label{subsec:obs-vs-decidability}

\begin{lemma}[Observability does not imply decidability]
The existence of a non-trivial observable projection
$\Omega_{\mathrm{obs}}$ does not, by itself, guarantee decidability of
any non-trivial state predicate.
\end{lemma}

\noindent
This lemma formalizes the strict separation between observability
(the availability of observable signatures) and decidability (logical
invariance of predicates under observational equivalence). The remainder
of the paper relies on this separation as a structural constraint, not
as a methodological choice.
